\chapter{Thực nghiệm và phân tích kết quả}

\section{Mục tiêu và thiết lập thực nghiệm}

Ba bộ dữ liệu được dùng với ba vai trò khác nhau. Letter Recognition là bộ dữ liệu chính
để đánh giá cây quyết định trên đặc trưng đã được trích xuất; Handwritten Digits kiểm tra
khả năng áp dụng trên biểu diễn điểm ảnh; Covertype kiểm tra khả năng mở rộng trên dữ liệu
lớn. Cách tổ chức này giúp tránh kết luận dựa trên một dạng dữ liệu duy nhất.

Các notebook được chạy trên Kaggle với Linux, 4 CPU logic và không bật GPU. Mô hình
\texttt{DecisionTreeClassifier} của scikit-learn huấn luyện và suy luận trên CPU. Mọi thí
nghiệm dùng \texttt{random\_state=42}, \texttt{test\_size=0.20} và phân tầng theo nhãn.
Không dùng tập test để chọn mô hình hoặc siêu tham số.

\begin{table}[H]
    \centering
    \caption{Vai trò và quy mô các bộ dữ liệu}
    \label{tab:dataset-summary}
    \small
    \begin{tabularx}{\textwidth}{@{}lXrrr@{}}
        \toprule
        \textbf{Bộ dữ liệu} & \textbf{Vai trò} & \textbf{Số mẫu} & \textbf{Đặc trưng} & \textbf{Số lớp} \\
        \midrule
        Letter Recognition & Bộ dữ liệu chính; đã loại trùng chính xác & 18.668 & 16 & 26 \\
        Handwritten Digits & So sánh chéo trên điểm ảnh $8\times8$ & 1.797 & 64 & 10 \\
        Covertype & Kiểm tra khả năng mở rộng & 581.012 & 54 & 7 \\
        \bottomrule
    \end{tabularx}
\end{table}

Letter Recognition được loại 1.332 dòng trùng chính xác trước khi chia dữ liệu, còn 18.668
mẫu. Các giá trị bị quy tắc IQR đánh dấu vẫn được giữ vì chúng nằm trong miền hợp lệ 0--15
của bộ dữ liệu. Decision Tree không yêu cầu chuẩn hóa thang đo hay one-hot encoding bổ sung
cho ba bộ dữ liệu đã được chuẩn bị.

\section{Chỉ số đánh giá}

Kết quả được đánh giá bằng accuracy, error rate và macro-F1. Macro-F1 được ưu tiên khi
phân bố lớp không cân bằng vì mỗi lớp đóng góp như nhau vào trung bình. Độ chênh giữa
accuracy trên train và test được dùng như dấu hiệu overfitting. Độ sâu và số lá phản ánh
độ phức tạp, trong khi thời gian nạp dữ liệu, huấn luyện, dự đoán và toàn pipeline phản ánh
chi phí tính toán.

Với Covertype, hai cấu hình được so sánh: cây baseline dùng Gini không giới hạn độ sâu và
cây regularized với \texttt{max\_depth=20}, \texttt{min\_samples\_leaf=5}. Việc chọn mô
hình dựa trên macro-F1 trung bình của 5-fold cross-validation chỉ trên tập train. Kết quả CV
được trình bày dưới dạng trung bình $\pm$ độ lệch chuẩn.

\section{Kết quả tổng hợp trên tập test}

\begin{table}[H]
    \centering
    \caption{Kết quả baseline trên tập test Kaggle}
    \label{tab:baseline-results}
    \small
    \begin{tabular}{@{}lrrrrr@{}}
        \toprule
        \textbf{Bộ dữ liệu} & \textbf{Accuracy} & \textbf{Macro-F1} & \textbf{Error} & \textbf{Độ sâu} & \textbf{Số lá} \\
        \midrule
        Letter Recognition & 0,8677 & 0,8677 & 0,1323 & 31 & 1.965 \\
        Handwritten Digits & 0,8250 & 0,8230 & 0,1750 & 14 & 142 \\
        Covertype & \textbf{0,9389} & \textbf{0,9034} & \textbf{0,0611} & 41 & 23.956 \\
        \bottomrule
    \end{tabular}
\end{table}

Ba cây baseline đều đạt accuracy train bằng 1,0. Điều này cho thấy cây không cắt tỉa có
khả năng ghi nhớ hoàn toàn tập huấn luyện. Khoảng cách train-test lần lượt là 0,1323 trên
Letter Recognition, 0,1750 trên Digits và 0,0611 trên Covertype. Vì vậy, kết quả test phải
được đọc cùng độ phức tạp của cây thay vì chỉ nhìn vào accuracy train.

\subsection{Letter Recognition}

Baseline đạt accuracy 0,8677 và macro-F1 0,8677 trên 3.734 mẫu test. Hai giá trị gần nhau
cho thấy chất lượng tương đối đồng đều giữa 26 lớp. Tuy nhiên, độ sâu 31 và 1.965 lá cùng
accuracy train 1,0 cho thấy mô hình vẫn overfit; đây là cơ sở để thử pruning hoặc điều chỉnh
\texttt{max\_depth}, \texttt{min\_samples\_leaf} ở giai đoạn cải tiến.

\begin{figure}[H]
    \centering
    \includegraphics[width=0.78\textwidth]{dt_letter_baseline__confusion_matrix.png}
    \caption{Ma trận nhầm lẫn của Decision Tree baseline trên Letter Recognition.}
    \label{fig:letter-confusion}
\end{figure}

\begin{figure}[H]
    \centering
    \includegraphics[width=0.92\textwidth]{dt_letter_baseline__feature_importance.png}
    \caption{Mức quan trọng của các đặc trưng trên Letter Recognition.}
    \label{fig:letter-importance}
\end{figure}

\subsection{Handwritten Digits}

Digits đạt accuracy 0,8250 và macro-F1 0,8230. Đây là kết quả thấp nhất trong ba bộ dữ
liệu và có train-test gap lớn nhất. Một cây đơn lẻ phải chia không gian 64 chiều dựa trên
từng điểm ảnh, nên nhạy với biến thiên nét viết và không khai thác được cấu trúc không gian
tốt như các mô hình chuyên cho ảnh.

\begin{figure}[H]
    \centering
    \includegraphics[width=0.76\textwidth]{dt_digits_baseline__confusion_matrix.png}
    \caption{Ma trận nhầm lẫn của Decision Tree baseline trên Handwritten Digits.}
    \label{fig:digits-confusion}
\end{figure}

\begin{figure}[H]
    \centering
    \includegraphics[width=0.98\textwidth]{dt_digits_baseline__misclassified_examples.png}
    \caption{Một số chữ số bị phân loại sai trên tập test.}
    \label{fig:digits-errors}
\end{figure}

\subsection{Covertype và đánh đổi hiệu năng - độ phức tạp}

Covertype có phân bố lớp lệch mạnh: lớp 1 và 2 chiếm đa số, còn lớp 4 có ít mẫu nhất. Vì
vậy macro-F1 0,9034 được dùng cùng accuracy 0,9389 để tránh đánh giá quá lạc quan dựa trên
hai lớp lớn. Mô hình baseline được chọn bởi 5-fold CV và đạt kết quả test tốt nhất, nhưng
cây sâu 41 với 23.956 lá.

\begin{figure}[H]
    \centering
    \includegraphics[width=0.95\textwidth]{dt_covertype_scalability__class_distribution.png}
    \caption{Phân bố bảy lớp của bộ dữ liệu Covertype.}
    \label{fig:covertype-distribution}
\end{figure}

\begin{table}[H]
    \centering
    \caption{So sánh hold-out giữa hai cấu hình trên Covertype}
    \label{tab:covertype-holdout}
    \small
    \begin{tabular}{@{}lrrrrr@{}}
        \toprule
        \textbf{Mô hình} & \textbf{Train acc.} & \textbf{Test acc.} & \textbf{Macro-F1} & \textbf{Độ sâu} & \textbf{Số lá} \\
        \midrule
        Baseline & 1,0000 & \textbf{0,9389} & \textbf{0,9034} & 41 & 23.956 \\
        Regularized & 0,9268 & 0,8988 & 0,8432 & \textbf{20} & \textbf{9.679} \\
        \bottomrule
    \end{tabular}
\end{table}

Cấu hình regularized giảm 59,6\% số lá và giảm train-test gap từ 0,0611 xuống 0,0279,
nhưng accuracy test giảm 0,0400 và macro-F1 giảm 0,0602. Với hai cấu hình hiện tại,
baseline vẫn là lựa chọn tốt hơn theo macro-F1; regularized minh họa rõ cái giá phải trả khi
giảm mạnh độ phức tạp. Các thí nghiệm tiếp theo nên tìm vùng trung gian thay vì kết luận
rằng mọi hình thức regularization đều làm mô hình kém hơn.

\begin{figure}[H]
    \centering
    \includegraphics[width=\textwidth]{dt_covertype_scalability__model_comparison.png}
    \caption{Đánh đổi giữa hiệu năng hold-out và độ phức tạp trên Covertype.}
    \label{fig:covertype-comparison}
\end{figure}

\begin{figure}[H]
    \centering
    \includegraphics[width=0.78\textwidth]{dt_covertype_scalability__confusion_matrix.png}
    \caption{Ma trận nhầm lẫn của mô hình Covertype được chọn bằng cross-validation.}
    \label{fig:covertype-confusion}
\end{figure}

\begin{figure}[H]
    \centering
    \includegraphics[width=0.94\textwidth]{dt_covertype_scalability__feature_importance.png}
    \caption{Hai mươi đặc trưng quan trọng nhất của cây Covertype baseline.}
    \label{fig:covertype-importance}
\end{figure}

\section{Độ ổn định qua cross-validation}

\begin{table}[H]
    \centering
    \caption{Kết quả 5-fold cross-validation trên tập train Covertype}
    \label{tab:covertype-cv}
    \small
    \begin{tabular}{@{}lccc@{}}
        \toprule
        \textbf{Mô hình} & \textbf{Accuracy} & \textbf{Macro-F1} & \textbf{Fit time (giây)} \\
        \midrule
        Baseline & \textbf{0,9330 $\pm$ 0,0011} & \textbf{0,8923 $\pm$ 0,0011} & 8,6710 $\pm$ 0,0852 \\
        Regularized & 0,8966 $\pm$ 0,0019 & 0,8388 $\pm$ 0,0028 & \textbf{7,4213 $\pm$ 0,1392} \\
        \bottomrule
    \end{tabular}
\end{table}

Độ lệch chuẩn nhỏ cho thấy kết quả ổn định giữa các fold. Baseline cao hơn regularized
0,0364 điểm accuracy và 0,0535 điểm macro-F1, trong khi thời gian fit trung bình chỉ tăng
khoảng 1,25 giây mỗi fold. Do đó, tiêu chí macro-F1 chọn baseline là hợp lý trong phạm vi
hai cấu hình đã thử.

\section{Thời gian chạy trên Kaggle}

\begin{table}[H]
    \centering
    \caption{Thời gian của các pipeline Kaggle, đơn vị giây}
    \label{tab:runtime}
    \small
    \begin{tabular}{@{}lrrrr@{}}
        \toprule
        \textbf{Bộ dữ liệu} & \textbf{Nạp dữ liệu} & \textbf{Huấn luyện} & \textbf{Dự đoán} & \textbf{Toàn pipeline} \\
        \midrule
        Letter Recognition & 0,0419 & 0,1264 & 0,0088 & 6,3916 \\
        Handwritten Digits & 0,0424 & 0,0239 & 0,0030 & 5,1282 \\
        Covertype & 86,7835 & 10,5151 & 0,2834 & 162,0717 \\
        \bottomrule
    \end{tabular}
\end{table}

Hai bộ dữ liệu nhỏ hoàn tất trong vài giây. Với Covertype, thời gian nạp dữ liệu chiếm
53,6\% toàn pipeline và lớn hơn thời gian huấn luyện cây. Tổng thời gian còn bao gồm hai
mô hình hold-out, 5-fold CV, sinh biểu đồ và đóng gói artifact. Điều này cho thấy nút thắt
trong notebook hiện tại không chỉ nằm ở thuật toán Decision Tree mà còn ở I/O và quy trình
đánh giá. GPU không giúp cải thiện bước huấn luyện vì estimator này chạy trên CPU.

\section{Kết luận thực nghiệm}

Kết quả từ ba bộ dữ liệu xác nhận Decision Tree là baseline nhanh, dễ diễn giải và có thể
chạy trên dữ liệu hơn nửa triệu mẫu. Tuy nhiên, accuracy train bằng 1,0 trên cả ba baseline
cho thấy nguy cơ overfitting rõ ràng. Covertype đạt hiệu năng test cao nhất và ổn định qua
5-fold CV, nhưng cây rất lớn. Regularization mạnh làm cây gọn hơn đáng kể nhưng giảm chất
lượng dự đoán. Vì vậy, bước cải tiến nên tìm mức cắt tỉa trung gian bằng cross-validation,
đồng thời tiếp tục theo dõi macro-F1, train-test gap, số lá và thời gian thay vì tối ưu duy
nhất accuracy.
